\documentclass[11pt]{amsart}
\usepackage[margin=1.2in]{geometry}
\usepackage{amsmath,amssymb,amsthm,mathtools,amscd}
\usepackage{mathrsfs}
\usepackage{booktabs}
\usepackage{xcolor}
\usepackage[colorlinks=true,linkcolor=blue!60!black,citecolor=blue!60!black,urlcolor=blue!60!black]{hyperref}
\usepackage{xurl}   % allow URL line breaks anywhere (TL >= 2018)
\emergencystretch=1.5em   % release stubborn overfull boxes

% ------------------------------------------------------------------ macros
\newcommand{\QQ}{\mathbb{Q}}
\newcommand{\ZZ}{\mathbb{Z}}
\newcommand{\CC}{\mathbb{C}}
\newcommand{\PP}{\mathbb{P}}
\newcommand{\GQ}{G_{\QQ}}
\newcommand{\GK}{G_{K}}
\newcommand{\Ni}{\operatorname{Ni}}
\newcommand{\Mtw}{M_{23}}
\newcommand{\sqm}{\sqrt{-23}}
\newcommand{\kap}{\kappa}
\newcommand{\Sym}{\operatorname{Sym}}
% AUTO-GENERATED by gen_results_tex.py -- do not edit
\newcommand{\RhoVal}{0.9812834897968066}
\newcommand{\NVal}{12187}
\newcommand{\QVal}{12207}
\newcommand{\LabelSixTable}{%
\begin{tabular}{ll}
\toprule
invariant & value in $K=\QQ(\sqm)$ \\
\midrule
$J1$ & $\dfrac{-4152+693\sqm}{10237}$ \\[2pt]
$J2$ & $\dfrac{-148227-34830\sqm}{142129}$ \\[2pt]
$J3$ & $\dfrac{-32137795-2206259\sqm}{18856554}$ \\[2pt]
$J4$ & $\dfrac{-44137861-681405\sqm}{18856554}$ \\[2pt]
$J5$ & $\dfrac{113993-1027623\sqm}{694434}$ \\[2pt]
$J6$ & $\dfrac{183063721+14235727\sqm}{182131716}$ \\[2pt]
\bottomrule
\end{tabular}
}
\newcommand{\GeneratorsTable}{%
\begin{tabular}{ll}
\toprule
$g_1$ ($2A$) & \texttt{(1,11)(2,23)(3,8)(4,16)(5,21)(7,20)(15,19)(18,22)} \\
$g_2$ ($23A$) & \parbox[t]{4.4in}{\raggedright\texttt{(1,2,11,10,16,9,6,3,23,19,20,14,21,17,4,8,22,5,18,15,13,7,12)}} \\[6pt]
$g_3$ ($23B$) & \parbox[t]{4.4in}{\raggedright\texttt{(1,2,3,4,10,11,12,7,19,18,8,6,9,16,17,21,22,5,14,20,13,15,23)}} \\
\bottomrule
\end{tabular}}
\newcommand{\ParamsTable}{%
\begin{tabular}{ll}
\toprule
working precision & $100$ decimal digits \\
$\rho$ (fundamental-domain radius) & $\RhoVal\ldots$ \\
truncation order $N$ & $\NVal$ \\
collocation points $Q$ & $\QVal$ \\
residual gate & $10^{-88}$ \\
label-6 residuals & $\le 10^{-91}$ (all four vectors) \\
low/high cross-check & $9.9\times10^{-14}$ \\
\bottomrule
\end{tabular}}
\newcommand{\RepsTable}{%
\begin{tabular}{ll}
\toprule
label $i$ & anchored $y_i$ (so the pair is $(g_1,y_i)$) \\
\midrule
$1$ & \parbox[t]{4.6in}{\raggedright\scriptsize\texttt{(1,4,19,10,20,12,9,13,6,3,16,22,15,21,14,8,2,17,7,5,11,18,23)}} \\[4pt]
$2$ & \parbox[t]{4.6in}{\raggedright\scriptsize\texttt{(1,23,21,19,8,6,16,9,13,12,14,20,3,2,15,18,10,4,11,17,7,5,22)}} \\[4pt]
$3$ & \parbox[t]{4.6in}{\raggedright\scriptsize\texttt{(1,4,21,12,8,7,5,19,2,20,6,9,18,13,14,3,16,17,10,22,15,11,23)}} \\[4pt]
$4$ & \parbox[t]{4.6in}{\raggedright\scriptsize\texttt{(1,20,17,21,12,2,13,4,11,10,23,6,9,15,22,19,18,16,5,8,7,3,14)}} \\[4pt]
$5$ & \parbox[t]{4.6in}{\raggedright\scriptsize\texttt{(1,17,15,4,18,3,21,9,8,20,13,10,22,12,19,23,6,11,5,7,14,16,2)}} \\[4pt]
$6$ & \parbox[t]{4.6in}{\raggedright\scriptsize\texttt{(1,6,18,16,22,9,3,14,13,2,4,19,11,10,21,17,8,23,20,5,7,15,12)}} \\[4pt]
$7$ & \parbox[t]{4.6in}{\raggedright\scriptsize\texttt{(1,22,8,12,9,21,11,3,13,15,5,14,23,19,7,10,4,17,2,20,16,18,6)}} \\[4pt]
\bottomrule
\end{tabular}}
\newcommand{\KappaCertsTable}{%
\begin{tabular}{ll}
\toprule
$i$ & $c_i$ with $c_i^{-1}(g_1y_i)c_i = y_{\sigma(i)}$ \\
\midrule
$1$ & \parbox[t]{4.6in}{\raggedright\scriptsize\texttt{(1,5,4,7,3,23)(2,11,21,16,20,8)(6,17,14)(9,12,13)(15,22)(18,19)}} \\[4pt]
$2$ & \parbox[t]{4.6in}{\raggedright\scriptsize\texttt{(1,23,3,7,4,5)(2,8,20,16,21,11)(6,14,17)(9,13,12)(15,22)(18,19)}} \\[4pt]
$3$ & \parbox[t]{4.6in}{\raggedright\scriptsize\texttt{(1,2)(3,8)(4,7)(6,17)(9,10)(11,23)(15,19)(16,20)}} \\[4pt]
$4$ & \parbox[t]{4.6in}{\raggedright\scriptsize\texttt{(1,5)(3,20)(4,16)(7,8)(10,12)(11,21)(14,17)(15,19)}} \\[4pt]
$5$ & \parbox[t]{4.6in}{\raggedright\scriptsize\texttt{(1,3,7,23,15,4)(2,19,16,11,8,20)(5,22)(6,9,10)(12,13,14)(18,21)}} \\[4pt]
$6$ & \parbox[t]{4.6in}{\raggedright\scriptsize\texttt{(1,11)(2,23)(4,7)(5,18)(6,10)(9,17)(16,20)(21,22)}} \\[4pt]
$7$ & \parbox[t]{4.6in}{\raggedright\scriptsize\texttt{(1,4,15,23,7,3)(2,20,8,11,16,19)(5,22)(6,10,9)(12,14,13)(18,21)}} \\[4pt]
\bottomrule
\end{tabular}}
\newcommand{\CGensLine}{$C=\langle h_1,h_2\rangle$, $h_1=$ \texttt{(1,8,5,7)(2,23)(3,21,20,11)(4,15,16,19)(9,13)(10,14,12,17)}, $h_2=$ \texttt{(1,4,3,22)(2,5,19,7)(8,18,11,16)(9,13)(14,17)(15,20,23,21)}}
\newcommand{\SexticTable}{%
\begin{tabular}{ll}
\toprule
\multicolumn{2}{l}{$f(T) = T^6 + \tfrac{1}{50019668646075143}\sum_k c_k T^k$ with $c_k$:} \\
\midrule
$T^{0}$ & $-139453053861313073$ \\
$T^{1}$ & $-868575878997768238$ \\
$T^{2}$ & $-759255603013669319$ \\
$T^{3}$ & $-55027953306246948$ \\
$T^{4}$ & $353749018385880145$ \\
$T^{5}$ & $206305139624557234$ \\
\bottomrule
\end{tabular}}
\newcommand{\IsoBetaTable}{%
\begin{tabular}{ll}
\toprule
$k$ & $r_k$ (numerator\,/\,denominator) \\
\midrule
$r_{0}$ & \parbox[t]{4.4in}{\raggedright\tiny\texttt{1139031854354334030139004760292567736326223637201380\allowbreak 9824818045843636895813575653251711824548075044675}\allowbreak\,/\,\allowbreak\texttt{2876099619276190692130997435382787257082398155096382\allowbreak 3324665343304091484039300414230962947781405761536}} \\[4pt]
$r_{1}$ & \parbox[t]{4.4in}{\raggedright\tiny\texttt{-700555257420276207380299130634698877119650508529204\allowbreak 36382040680840049664033269663315666076259606117145}\allowbreak\,/\,\allowbreak\texttt{2876099619276190692130997435382787257082398155096382\allowbreak 3324665343304091484039300414230962947781405761536}} \\[4pt]
$r_{2}$ & \parbox[t]{4.4in}{\raggedright\tiny\texttt{-631389859765933139842337994224873405676714039836722\allowbreak 49772219471359832896209510086913254392530340893153}\allowbreak\,/\,\allowbreak\texttt{1438049809638095346065498717691393628541199077548191\allowbreak 1662332671652045742019650207115481473890702880768}} \\[4pt]
$r_{3}$ & \parbox[t]{4.4in}{\raggedright\tiny\texttt{1285931227362825893728832979551508223464099626737486\allowbreak 8578492980729550747827078761297986286246826959095}\allowbreak\,/\,\allowbreak\texttt{1438049809638095346065498717691393628541199077548191\allowbreak 1662332671652045742019650207115481473890702880768}} \\[4pt]
$r_{4}$ & \parbox[t]{4.4in}{\raggedright\tiny\texttt{4216230577851849077693844100888088127968536559500513\allowbreak 2800860605265021844679263239610512790264358368879}\allowbreak\,/\,\allowbreak\texttt{2876099619276190692130997435382787257082398155096382\allowbreak 3324665343304091484039300414230962947781405761536}} \\[4pt]
$r_{5}$ & \parbox[t]{4.4in}{\raggedright\tiny\texttt{2750428657631107262887869693920120585531771976248552\allowbreak 39444109321284930781088274243147228225852595341}\allowbreak\,/\,\allowbreak\texttt{4050844534191817876240841458285615855045631204361101\allowbreak 87671342863437908225905639637055816165935292416}} \\[4pt]
\bottomrule
\end{tabular}}
\newcommand{\SexticTwoTable}{%
\begin{tabular}{ll}
\toprule
\multicolumn{2}{l}{$f_2(T) = T^6 + \tfrac{1}{141385642915876396665810361}\sum_k c_k T^k$ with $c_k$:} \\
\midrule
$T^{0}$ & $414128288417144984367919381$ \\
$T^{1}$ & $-3014801889733613820545273694$ \\
$T^{2}$ & $8281435636737711157987163503$ \\
$T^{3}$ & $-10236272033398981465535697652$ \\
$T^{4}$ & $5060936190487869829531262115$ \\
$T^{5}$ & $-489750524557256227855164622$ \\
\bottomrule
\end{tabular}}
\newcommand{\PtwoHundredLine}{\relax}
\newcommand{\PtwoHundredArtifactLine}{A full replication of the pipeline at $200$-digit working precision is in progress; its receipts will be added to the repository.}
\newcommand{\ConfirmationPhrase}{independently replicated by two descended coordinates with $\sim70$ digits of confirmation margin}
   % machine-generated numerical results (gen_results_tex.py)

\theoremstyle{plain}
\newtheorem{theorem}{Theorem}[section]
\newtheorem{numresult}[theorem]{Numerical Result}
\newtheorem{lemma}[theorem]{Lemma}
\newtheorem{proposition}[theorem]{Proposition}
\newtheorem{corollary}[theorem]{Corollary}
\theoremstyle{definition}
\newtheorem{definition}[theorem]{Definition}
\newtheorem{remark}[theorem]{Remark}
\newtheorem{convention}[theorem]{Convention}

\title[The Galois action on the seven $\Mtw$ Hurwitz points]{The
  Galois action on the seven $M_{23}$ Hurwitz points}
\author{Alexander West Templeton}
\address{Independent researcher}
\email{alex@templeton.cloud}
\date{\today}

\begin{document}

\begin{abstract}
Huang, Jackson, Lee, Poonen, Pries, and Zhang recently realized
$\Mtw$ as a Galois group over $\QQ$ by finding a $\GQ$-fixed point
among the seven inner Nielsen classes of type $(2A,23A,23B)$, and
noted that they had no conceptual explanation for this fixed point.
We determine the Galois action on the full seven-point Hurwitz
fiber.
%
The odd cardinality of this $\GQ$-stable fiber already forces a real
Hurwitz point, since complex conjugation is an involution. We compute
its action exactly. With
branch points $(0,\sqm,-\sqm)$, it acts on the seven classes as
$(1\,2)(5\,7)$, fixing $3$, $4$, $6$. Thus there are exactly three
real Hurwitz points; the existence of a $\QQ$-point is a further
arithmetic fact. We then compute the associated
genus-$4$ quotient curves by numerical uniformization and construct
marked Petri invariants which we prove algebraic and
$\GK$-equivariant for $K=\QQ(\sqm)$. Trace and norm along the
marking descend these to algebraic $\GQ$-equivariant invariants of
the unmarked fiber.
%
Integer-relation recognition, \ConfirmationPhrase, gives a
high-precision numerical determination of the arithmetic
decomposition
\[
  \{1,\dots,7\}=\{6\}\sqcup\{1,2,3,4,5,7\}.
\]
The two coordinates generate the same sextic residue field $F$; it
has signature $(2,2)$, and its Galois closure has group $S_6$. Thus,
subject only to the stated integer-relation recognitions,
\[
  \rho(\GQ)=\operatorname{Stab}_{S_7}(6)\cong S_6,
  \qquad
  \rho(\GK)=S_6.
\]
Arithmetic fixes exactly one point and acts maximally on the
remaining six.
%
This symmetry breaking is not forced by the finite-level
Grothendieck--Teichm\"uller constraints: using Guillot's description
of $\mathcal{GT}_1(G)$ for finite simple groups, we prove that
$\mathcal{GT}_1(\Mtw)$ induces the full $S_7$ on these seven Nielsen
classes. Finally, under the same recognition hypothesis, at the
prime $23$ the seven-point Hurwitz algebra decomposes locally as
$\QQ_{23}\times F_2\times F_4$, where $F_2$ and $F_4$ are totally
ramified tame extensions of degrees $2$ and $4$. Hence inertia has
cycle type $(1)(2)(4)$, making label $6$ the unique inertia-fixed
point at $23$ and giving a concrete local target for a conceptual
explanation of the exceptional rational point.
\end{abstract}

\maketitle

% ================================================================ INTRO
\section{Introduction}

The inverse Galois problem for the sporadic simple group $\Mtw$ resisted
solution longer than for any other sporadic group. It was settled in
2026 by Huang, Jackson, Lee, Poonen, Pries, and Zhang \cite{HJLPPZ},
who studied the Hurwitz space of Galois covers of $\PP^1$ with
monodromy group $\Mtw$, ramification type $(2A,23A,23B)$, and branch
points $(0,\sqm,-\sqm)$. The associated inner Nielsen class
\[
  \Ni(2A,23A,23B)^{\mathrm{in}}
  \;=\;
  \bigl\{(x,y,z)\in 2A\times 23A\times 23B :
     xyz=1,\ \langle x,y\rangle=\Mtw\bigr\}\big/\Mtw
\]
consists of exactly seven simultaneous-conjugacy classes; the Hurwitz
space is accordingly a degree-$7$ cover of the configuration base, its
fiber above $(0,\sqm,-\sqm)$ consisting of seven points permuted by
$\GQ=\operatorname{Gal}(\overline{\QQ}/\QQ)$. The authors of
\cite{HJLPPZ} located a $\GQ$-fixed point among the seven --- yielding a
regular realization of $\Mtw$ over $\QQ$ --- and remarked that they
``do not have a conceptual explanation'' for its existence.

This note contributes two computations toward that explanation.

\subsection*{Complex conjugation}
Before computing the permutation, parity already gives a real point
(Corollary~\ref{cor:parity}): an involution on seven points can pair
off only an even number of them. This forces at least one real
embedding of a field of moduli, but neither the number of real
points nor a $\GQ$-fixed point.

First, we determine unconditionally how the subgroup generated by
complex conjugation acts on the seven points
(Theorem~\ref{thm:kappa}): in our labeling
(Convention~\ref{conv:group}) it is the involution
\[
  \kap \;=\; (1\,2)(5\,7),
\]
fixing the three points $3,4,6$. In particular exactly three of the
seven points are fixed by the chosen complex conjugation (so their
residue fields admit the corresponding real embedding), and the
$\GQ$-orbit partition of the seven refines to one of the eight shapes
listed in Corollary~\ref{cor:menu}. The proof combines the branch cycle argument
of Fried \cite{Fried77, DebesDouai} with a finite certificate: after
anchoring the $2A$-coordinate, complex conjugation acts on the
anchored slice by $y \mapsto x_0y$, and its effect on the seven
classes is witnessed by seven explicit conjugating elements
(\S\ref{sec:kappa}).
The three assertions are distinct: real existence is forced by
parity, the exact count of three is computed, and rationality
requires a point fixed by all of $\GQ$. The rational point of
\cite{HJLPPZ} supplies this additional arithmetic input.

\subsection*{The orbit partition}
Second, we compute the seven covers themselves. Each quotient curve
$X_i$ (the quotient of the Galois cover by a point stabilizer) is a
curve of genus $4$; we compute its canonical differentials and Petri
quadric by the power-series method of
Klug--Musty--Schiavone--Voight \cite{KMSV} at $100$ decimal
digits, using the standalone Arnoldi solver of \cite{BelyiRepo}
(\S\ref{sec:powser}). From the Petri quadric of the canonical curve we
extract \emph{marked invariants}: ratios of quadric coefficients that
are invariant both under the projective scale and under the one
remaining analytic normalization freedom, hence are algebraic numbers
attached to the marked cover, compatible with the Galois action
(\S\ref{sec:invariants}). For the $\GQ$-fixed point (label $6$) all
six invariants are recognized as elements of $K=\QQ(\sqm)$ with small
height --- for example
\[
  J_1 \;=\; \frac{-4152+693\sqm}{10237},
\]
with $\sim 85$ digits of confirmation margin --- validating the method
against known ground truth. Applying the same procedure to all seven
points and factoring the resulting degree-$7$ polynomials over $K$
yields the $\GK$-orbit partition, which combined with $\kap$
determines the $\GQ$-partition: it is $1+6$, and the permutation
image is the maximal group compatible with one fixed point
(Numerical Result~\ref{thm:partition}).

\subsection*{Verification}
Every computational claim in this note is backed by a machine-readable
receipt (\S\ref{sec:artifacts}); the census and the $\kap$-action are
verified by three independent implementations, including a Magma
referee run; a Lean~4 formalization of the finite combinatorial core
(kernel-checked \texttt{decide} certificates with a single
quarantined \texttt{native\_decide} enumeration) is in progress and
is not claimed as part of the verification here. The analytic input --- that complex
conjugation acts through the branch cycle argument --- enters only
through the published results cited above.

\subsection*{Relation to \cite{HJLPPZ}}
The authors of \cite{HJLPPZ} have announced a revision containing more
details, including the action of complex conjugation, which they
computed independently. The present computations were performed
without access to that revision; the labeled representatives in
Appendix~\ref{app:reps} make the two computations directly comparable.

% ============================================================== CENSUS
\section{The seven Nielsen classes}\label{sec:census}

\begin{convention}\label{conv:group}
$\Mtw$ denotes the copy $\langle g_1,g_2\rangle\subset S_{23}$
generated by the two permutations of \cite{HJLPPZ},
reproduced in Appendix~\ref{app:reps}; $g_1\in 2A$, $g_2\in 23A$,
$g_3:=(g_1g_2)^{-1}\in 23B$. Products act on the right of the first
factor: $(pq)(i)=p(q(i))$. We write $C=C_{\Mtw}(g_1)$, of order
$2688$. We use our \emph{own} labels $1,\dots,7$ for the seven
Nielsen classes, fixed by the anchored representatives of
Appendix~\ref{app:reps} and chosen so that label $6$ is the class
of the generating triple printed in \cite{HJLPPZ} (which does not
publish a numbering); the appendix permits direct comparison with
any numbering in their announced revision.
\end{convention}

\begin{proposition}[census]\label{prop:census}
The set
\[
  T=\bigl\{(x,y)\in 2A\times 23A :\ (xy)^{-1}\in 23B,\
  \langle x,y\rangle = \Mtw\bigr\}
\]
carries exactly seven orbits under simultaneous conjugation, with
representatives $(g_1,y_1)$, \dots, $(g_1,y_7)$ listed in
Appendix~\ref{app:reps} (numbered as in
Convention~\ref{conv:group}). Moreover:
\begin{enumerate}
\item every orbit meets the anchored slice $\{(g_1,y)\}$, and two
  anchored pairs are conjugate iff they are $C$-conjugate;
\item the set $Y=\{y : (g_1,y)\in T\}$ has exactly
  $18816 = 7\cdot 2688$ elements, partitioned into seven
  \emph{regular} $C$-orbits (equivalently, $Z(\Mtw)=1$ and each pair
  has trivial simultaneous centralizer);
\item the generation condition is automatic: every $y\in 23A$ with
  $(g_1y)^{-1}\in 23B$ generates $\Mtw$ together with $g_1$.
\end{enumerate}
\end{proposition}

\begin{proof}[Proof (certificate)]
(1) is the orbit--stabilizer transport for the conjugation action on
the $2A$-coordinate. (2) and (3) are finite assertions verified by
exhaustive enumeration: the classes $23A$, $23B$ (each of size
$443\,520$) are materialized as conjugation orbits, $Y$ is filtered,
and the seven $C$-orbits of the $y_i$ are checked to be disjoint, of
size $2688$ each, and to cover $Y$; generation for the seven
representatives is certified by verified stabilizer chains
(order $10\,200\,960$ for each $\langle g_1,y_i\rangle$), and
generation is constant on orbits. The enumeration has been performed
by three independent implementations (\S\ref{sec:artifacts});
the in-progress Lean formalization packages the same certificates.
\end{proof}

\begin{remark}
The identity $|Y| = 7\cdot|C|$, i.e.\ the regularity of all seven
orbits, is equivalent to the vanishing of the simultaneous centralizer
of each generating pair, which holds since $Z(\Mtw)=1$. This is what
makes the anchored certificate small: seven orbit representatives,
two generators of $C$, and seven conjugators suffice.
\end{remark}

% =============================================================== KAPPA
\section{Complex conjugation: the involution \texorpdfstring{$\kap$}{kappa}}
\label{sec:kappa}

Write $H$ for the unmarked degree-$7$ Hurwitz fiber over $\QQ$ from
the Introduction. Its real points are precisely the geometric
points fixed by the chosen complex conjugation. This is the
$\GQ$-stable unmarked fiber; the $\sqm$-marked data used later are
naturally over $K=\QQ(\sqm)$.

Fix the branch points $(b_1,b_2,b_3)=(0,\sqm,-\sqm)$ and a real base
point, and order the standard generators of
$\pi_1(\PP^1\smallsetminus\{b_1,b_2,b_3\})$ so that the slots carry
the classes $(2A,23A,23B)$ as in \cite{HJLPPZ}. Complex conjugation
fixes $b_1$ and exchanges $b_2\leftrightarrow b_3$; the branch cycle
argument \cite{Fried77,DebesDouai} computes its effect on Nielsen
classes for a conjugation-symmetric bouquet:
\[
  \kap(x,y,z) \;=\; (x,\; z^{-1},\; y^{-1}).
\]
(For a triple with $xyz=1$ this is again a triple with product $1$ and
classes $(2A,23A,23B)$, since inversion exchanges $23A$ and
$23B$.\footnote{The generic form of the formula is
$(x^{-1},z^{-1},y^{-1})$; here $x$ is an involution.})

Reversal--inversion gives $(z^{-1},y^{-1},x^{-1})$, in class order
$(23A,23B,2A)$.
Cyclically rotating that tuple to restore the $2A$ slot gives the
displayed formula. The anchored certificates below use the restored
class order $(2A,23A,23B)$.

\begin{corollary}[parity: forced real existence]\label{cor:parity}
The seven-point Hurwitz fiber $H$ has a real point, independently of
the explicit computation of $\kap$.
\end{corollary}

\begin{proof}
Complex conjugation is an involution, so its orbits have sizes $1$
or $2$. Hence $\#H(\mathbb{R})+2m=7$ for some integer $m\ge0$;
the number of fixed points is odd and positive. This uses only the
seven-point census and stability under complex conjugation.
\end{proof}

\begin{theorem}\label{thm:kappa}
On the seven classes of Proposition~\ref{prop:census}, complex
conjugation acts by the involution
$\kap=(1\,2)(5\,7)$, fixing the labels $3$, $4$, $6$.
\end{theorem}

\begin{proof}[Proof (certificate)]
On the anchored slice the map takes the pleasant form
$y \mapsto g_1y$: indeed
$\kap(g_1,y,z)=(g_1,z^{-1},y^{-1})$ with
$z^{-1}=(g_1y)^{-1\,-1}=g_1y$. For each label $i$ an explicit
$c_i\in C$ with
\[
  c_i^{-1}\,(g_1 y_i)\,c_i \;=\; y_{\sigma(i)},\qquad
  \sigma=(1\,2)(5\,7),
\]
is exhibited in the artifact repository and re-verified by the Magma
referee; by Proposition~\ref{prop:census}(1) this determines the
action on classes. Involutivity of the induced $\sigma$ and the
membership certificates $c_i\in C\cap\Mtw$ are part of the same
receipt.
\end{proof}

\begin{remark}[forced, computed, arithmetic]
Corollary~\ref{cor:parity} forces $\#H(\mathbb{R})\ge1$, whereas
Theorem~\ref{thm:kappa} computes $\#H(\mathbb{R})=3$. Equivalently,
the degree-$7$ \'etale coordinate algebra of $H$ has exactly three
real embeddings, counted across its residue fields (fields of
moduli). A $\QQ$-point requires a geometric point fixed by every
element of $\GQ$, rather than by complex conjugation alone. Thus
the exact certificates restrict possible rational points to labels
$3,4,6$, but do not imply that any of them is rational. They do not
exclude a transitive seven-point orbit, whose residue field would
have signature $(3,2)$ (Corollary~\ref{cor:menu}).

Label $6$ is the $\GQ$-fixed point established independently by
\cite{HJLPPZ}; that $\kap(6)=6$ is a consistency check the
certificate passes, not an input. Its uniqueness and the remaining
sextic orbit are the stronger arithmetic conclusions of Numerical
Result~\ref{thm:partition}, subject to the stated integer-relation
recognitions.
\end{remark}

\begin{corollary}[partition menu]\label{cor:menu}
Every $\GQ$-orbit is a union of $\kap$-blocks
$\{1,2\}$, $\{5,7\}$, $\{3\}$, $\{4\}$, $\{6\}$.
\begin{enumerate}
\item From Theorem~\ref{thm:kappa} alone, the partition shape of the
  seven is one of the twelve shapes
  $2{+}2{+}1{+}1{+}1$, $2{+}2{+}2{+}1$, $3{+}2{+}2$,
  $3{+}2{+}1{+}1$, $3{+}3{+}1$, $4{+}1{+}1{+}1$, $4{+}2{+}1$,
  $4{+}3$, $5{+}1{+}1$, $5{+}2$, $6{+}1$, $7$.
\item Since label $6$ is $\GQ$-fixed \cite{HJLPPZ}, the orbit $\{6\}$
  splits off and the partition of the remaining six is one of the
  eight shapes
  \[
    2{+}2{+}1{+}1,\quad 2{+}2{+}2,\quad 3{+}2{+}1,\quad 3{+}3,\quad
    4{+}1{+}1,\quad 4{+}2,\quad 5{+}1,\quad 6 .
  \]
\end{enumerate}
Numerical Result~\ref{thm:partition} selects the true shape.
\end{corollary}

% ============================================================== POWSER
\section{Power series models of the seven quotient curves}
\label{sec:powser}

Let $\Delta=\Delta(23,2,23)$ be the hyperbolic triangle group with the
order-$23$ vertex uniformizing the branch point $\sqm$ (slot $23A$),
and for each label $i$ let $\Gamma_i\le\Delta$ be the index-$23$
subgroup determined by the permutation triple of $(g_1,y_i)$. The
quotient $X_i=\Gamma_i\backslash\mathcal{H}$ is the genus-$4$ quotient
curve of the $i$-th cover; by Riemann--Hurwitz
$g = 1+\tfrac{1}{2}(23-1-15-1)=4$, where $15$ is the number of cycles
of the $2A$-permutation.

We compute a basis of the weight-$2$ cusp forms
$S_2(\Gamma_i)=H^0(X_i,\Omega^1)$ as power series in the disc
uniformizer $w$ centered at the elliptic fixed point of the order-$23$
vertex, following \cite{KMSV,VoightWillis}: with
$\rho=\RhoVal\ldots$ (the radius of the fundamental domain),
truncation order $N=\NVal$, and $Q=\QVal$ collocation points on
$|w|=\rho$, the coefficient vector of each form is a fixed point of a
structured evaluation--interpolation operator $A$, computed by the
Arnoldi iteration implemented in the standalone solver of
\cite{BelyiRepo} at working precision $100$ digits.

Validation is layered:
\begin{enumerate}
\item the geometry stack (coset representatives, fundamental-domain
  reduction, automorphy factors) passes a battery of exact receipts,
  including the identity $\operatorname{Im}(\gamma z) =
  \operatorname{Im}(z)/|cz+d|^2$ at every collocation point and the
  radius bound $|w_{\mathrm{red}}|\le\rho$;
\item at low precision the operator $A$, rebuilt from the emitted
  input file alone, has numerical fixed space of dimension exactly
  $4=\dim S_2(\Gamma_i)$, with singular-value separation
  $10^{-15}$ vs.\ $10^{-1.1}$ --- an end-to-end structural check;
\item the $100$-digit solve validates each returned vector $x$ by
  the finite-operator residual $\|Ax-x\|/\|x\|$; for label $6$ all
  four residuals are $\le 10^{-91}$, against an acceptance gate of
  $10^{-88}$ (this validates the numerical operator, not the
  truncation-to-algebraic step, which is addressed in
  \S\ref{sec:invariants});
\item the $100$-digit coefficients agree with the independently
  computed low-precision fixed space to $\sim 10^{-13}$, the roundoff
  floor of the latter.
\end{enumerate}
For the remaining labels the validated residuals (item 3) are
$1.2\times10^{-91}$ (label $1$), $4.6\times10^{-92}$ ($2$),
$3.3\times10^{-92}$ ($3$), $8.9\times10^{-92}$ ($4$),
$1.7\times10^{-92}$ ($5$), and $2.3\times10^{-92}$ ($7$), all
against the same $10^{-88}$ gate.

% =========================================================== INVARIANTS
\section{Marked Petri invariants}\label{sec:invariants}

Write $f_0,f_1,f_2,f_3$ for the basis of $S_2(\Gamma_i)$ echelonized
with respect to the $w$-expansion,
$f_j = w^j + O(w^4)$. Since $X_i$ is not hyperelliptic (this is
checked, not assumed: see below), the canonical embedding
$X_i\hookrightarrow\PP^3$ realizes $X_i$ as the complete intersection
of a unique quadric $Q_i$ and a cubic modulo $Q_i$ (Petri). We compute
$Q_i=\sum_{j\le k} q_{jk}\,x_jx_k$ from the one-dimensional kernel of
the evaluation matrix of the ten products $f_jf_k$ --- the kernel
dimension being exactly $1$ is the non-hyperellipticity gate.

\begin{lemma}[forced filtration structure]\label{lem:filtration}
For every label, $q_{00}=q_{01}=0$, $q_{02}=-q_{11}$, and
$q_{03}=-q_{12}$ \textup{(}in the echelon basis, after scaling\textup{)}.
\end{lemma}

\begin{proof}
The relation $\sum q_{jk} f_jf_k=0$ has $w$-expansion beginning
\[
  q_{00} + q_{01}w + (q_{02}+q_{11})w^2 + (q_{03}+q_{12})w^3 +
  O(w^4),
\]
because $f_jf_k = w^{j+k}+O(w^4)$ and no products contribute below
order $4$ otherwise. Each displayed coefficient must vanish.
\end{proof}

Numerically the four vanishing combinations are observed at
$10^{-92}$, an additional receipt. The moduli-carrying data is thus
the vector $(q_{11},q_{12},q_{13},q_{22},q_{23},q_{33})$ up to the
following normalization freedom.

\begin{lemma}[weights]\label{lem:weights}
The echelon basis is unique given the local parameter $w$; the
residual freedom is $w\mapsto\lambda w$ \textup{(}$\lambda\in\CC^*$%
\textup{)}, under which the re-echelonized basis satisfies
$f_j\rightsquigarrow \lambda^{-j}f_j(\lambda\,\cdot)$ and the quadric
coefficients transform by
\[
  q_{jk} \;\longmapsto\; \mu\,\lambda^{j+k}\,q_{jk}
  \qquad(\mu\in\CC^* \text{ a common scale}).
\]
Consequently any ratio $\prod q_{j_tk_t}\big/\prod q_{l_tm_t}$ with
equally many factors up and down and matching total weight
$\sum(j_t+k_t)=\sum(l_t+m_t)$ is independent of all normalization
choices. \textup{(}The sign of the weight exponent depends on
whether one transforms functions or differentials; only weight
\emph{differences} matter for such ratios --- see
Convention~\ref{conv:weights} for the form-level convention used in
the algebraic argument.\textup{)}
\end{lemma}

We call such ratios \emph{marked invariants}. The next proposition
--- the hinge of the method --- shows they are algebraic and
$\GK$-equivariant. Its proof is purely algebraic: no numerical
statement enters. It rests on three observations: adapted parameters
form a $\mu_{23}$-torsor; the analytic parameter agrees with any
adapted one to order $>12$; and on a genus-$4$ curve the Petri
quadric is determined by $12$-jets, since $\deg 2K_X=12$. The
structural reason the argument closes is the inequality
\[
  \boxed{\,23 \;>\; \deg 2K_X \;=\; 4g-4 \;=\; 12\,.}
\]

\begin{definition}[adapted parameter]\label{def:adapted}
Let $(X,\phi)$ be a cover in the family over a finite extension
$L/K$ contained in $\overline{\QQ}$, and $P$ the unique (totally
ramified) point above $\sqm$. A formal parameter $t$ at $P$ is \emph{adapted} if
$\phi = \sqm + t^{23}$ exactly in
$\widehat{\mathcal{O}}_{X_{\overline L},P}$. Write
$\mathscr{T}_P=\{t \text{ adapted}\}$.
\end{definition}

\begin{lemma}[torsor]\label{lem:adapted}
The functor $\mathscr{T}_P$ is a $\mu_{23}$-torsor over $L$ --- any
two adapted parameters differ by $t\mapsto\mu t$, $\mu^{23}=1$ ---
and $\mathscr{T}_P(\overline{L})\neq\varnothing$. Equivalently:
adapted parameters exist after a finite extension of $L$, and need
not exist over $L$ itself. The torsor is functorial under
$\operatorname{Gal}(\overline{L}/L)$ for $L\supseteq K$: if $t$ is
adapted for $(X,\phi,P)$ then $t^{\sigma}$ is adapted for
$(X^{\sigma},\phi^{\sigma},P^{\sigma})$ \textup{(}using
$\sigma(\sqm)=\sqm$\textup{)}. Moreover, if $w$ is the analytic disc
parameter of \S\ref{sec:powser} \textup{(}canonical up to
$w\mapsto\lambda w$\textup{)}, then every adapted $t$ satisfies
\[
  t \;=\; \lambda\, w\,\bigl(1 + u_1w^{23} + u_2w^{46}+\cdots\bigr),
  \qquad
  dt \;=\; \lambda\,\bigl(1+O(w^{23})\bigr)\,dw
\]
for some $\lambda\in\CC^*$: the comparison unit is a power series in
$w^{23}$.
\end{lemma}

\begin{proof}
Writing $\phi-\sqm = c\,\tilde t^{\,23}\bigl(1+h(\tilde t)\bigr)$ in
any parameter $\tilde t$ over $L$, an adapted parameter is
$t=\tilde t\,\bigl(c(1+h)\bigr)^{1/23}$; the $23$rd root of the unit
$1+h$ exists in $L[[\tilde t\,]]$ (characteristic $0$), while the
root $c^{1/23}$ in general requires a finite extension --- whence
existence over $\overline L$ only, and the $\mu_{23}$-torsor
structure. Functoriality is immediate from
$\sigma(\phi-\sqm)=\phi^{\sigma}-\sqm$. For the analytic comparison:
$\phi-\sqm$ is invariant under the order-$23$ rotation
$w\mapsto\zeta_{23}w$ about the disc center, hence is a power series
in $w^{23}$; extracting the $23$rd root as above gives the displayed
form, and $dt$ follows by differentiation.
\end{proof}

\begin{convention}[weights on coefficients]\label{conv:weights}
Fix a strictly increasing vanishing sequence
$0\le a_0<a_1<a_2<a_3$ (the orders of vanishing at $P$ realized on
$H^0(X,\Omega^1)$; generically, and for all our labels by a per-label
numerical gate, $a_j=j$). The echelon basis w.r.t.\ a parameter $t$
is $\omega_j = \bigl(t^{a_j}+\sum_{n>a_j}a_{j,n}t^{\,n}\bigr)dt$
with $a_{j,a_k}=0$ for $k\neq j$. Under $t\mapsto\mu t$, after
renormalizing leading terms,
\[
  a_{j,n}\;\longmapsto\;\mu^{\,a_j-n}\,a_{j,n},
\]
and the Petri coefficients transform by
$q_{jk}\mapsto \tilde\mu\,\mu^{-(a_j+a_k)}q_{jk}$ for a common
scalar $\tilde\mu$; only weight \emph{differences} matter below, so
the sign convention in the exponent is immaterial for weight-matched
ratios. A ratio $\prod q_{j_tk_t}/\prod q_{l_tm_t}$ is
\emph{weight-$0$} if the factor counts agree and
$\sum(a_{j_t}+a_{k_t})=\sum(a_{l_t}+a_{m_t})$.
\end{convention}

\begin{lemma}[Petri from $12$-jets]\label{lem:jets}
Let $X$ be a non-hyperelliptic curve of genus $4$ and $P\in X$ any
point. The jet map
$j^{12}_P\colon H^0(X,2K_X)\to (2K_X)\otimes
\mathcal{O}_{X,P}/\mathfrak{m}_P^{13}$
is injective, and consequently the Petri line
$\ker\bigl(m\colon\Sym^2H^0(K_X)\to H^0(2K_X)\bigr)$ equals
$\ker\bigl(j^{12}_P\circ m\bigr)$: the quadric is determined exactly
by the coefficients of the products $\omega_j\omega_k$ through order
$12$ at $P$.
\end{lemma}

\begin{proof}
$\deg 2K_X = 4g-4 = 12$, so a nonzero global section of $2K_X$
vanishes at $P$ to order at most $12$; a section with vanishing
$12$-jet is therefore zero. The second statement is immediate.
\end{proof}

\begin{proposition}[equivariance]\label{prop:equivariance}
Let $J$ be any weight-$0$ ratio in the sense of
Convention~\ref{conv:weights}, defined on the locus where its
denominator factors are nonzero. Then:
\begin{enumerate}
\item $J(X,\phi)\in\overline{\QQ}$ is an invariant of the
  isomorphism class of the marked cover $(X,\phi,P)$, and for every
  $\sigma\in\GK$ the square
  \[
  \begin{array}{ccc}
    (X,\phi,P) & \xmapsto{\ \ J\ \ } & \overline{\QQ}\\[2pt]
    \big\downarrow{\scriptstyle\sigma} & & \big\downarrow{\scriptstyle\sigma}\\[2pt]
    (X^\sigma,\phi^\sigma,P^\sigma) & \xmapsto{\ \ J\ \ } & \overline{\QQ}
  \end{array}
  \]
  commutes;
\item the analytically computed value of \S\ref{sec:powser} equals
  $J(X,\phi)$.
\end{enumerate}
\end{proposition}

\begin{proof}
(1) Pull back to the torsor: for an adapted $t\in\mathscr{T}_P$
(over $\overline{L}$, Lemma~\ref{lem:adapted}) the flag of vanishing
orders, the echelon basis, the products
$\omega_j\omega_k\in H^0(2K_X)$, their $12$-jets at $P$ in the
parameter $t$, and the Petri line recovered from those jets
(Lemma~\ref{lem:jets}) are all algebraic constructions, compatible
with any field embedding; hence all $a_{j,n}$ and $q_{jk}$ lie in
$\overline{\QQ}$ and depend on the choice of $t$ only through the
$\mu_{23}$-action of Convention~\ref{conv:weights}. A weight-$0$
ratio is $\mu_{23}$-invariant, so $J$ descends from the torsor to
the marked cover itself. Functoriality of $\mathscr{T}_P$ under
$\GK$ (Lemma~\ref{lem:adapted}) then gives commutativity: $\sigma$
carries the entire diagram
\[
  \begin{CD}
    \Sym^2 H^0(K_X) @>{m}>> H^0(2K_X) @>{j^{12}_P}>>
      J^{12}_P(2K_X)\\
    @V{\sigma}VV @V{\sigma}VV @V{\sigma}VV\\
    \Sym^2 H^0(K_{X^\sigma}) @>{m^\sigma}>> H^0(2K_{X^\sigma})
      @>{j^{12}_{P^\sigma}}>> J^{12}_{P^\sigma}(2K_{X^\sigma})
  \end{CD}
\]
to the corresponding diagram for $(X^{\sigma},\phi^{\sigma})$, with
the right-hand horizontal maps injective, and the Petri line
functorially recovered on each row.

(2) By Lemma~\ref{lem:adapted}, $t=\lambda w(1+O(w^{23}))$ and
$dt=\lambda(1+O(w^{23}))dw$, so the transition between the echelon
computations in $w$ and in $t$ differs from the diagonal scaling
$\operatorname{diag}(\lambda^{a_j})$ only at jet orders $\ge 23$.
Since $23>12$, the two computations have \emph{identical} $12$-jets
of the products $\omega_j\omega_k$ up to that diagonal scaling and a
common scalar; by Lemma~\ref{lem:jets} the Petri lines agree up to
the same scaling, and weight-$0$ ratios kill it. No truncation
hypothesis is needed: the entire $2K$-linear system is blind to the
nonlinear discrepancy between $w$ and $t$, precisely because
$23>\deg 2K_X$.
\end{proof}

\begin{remark}[numerical cross-check]
Although the proof above is exact, the pipeline still runs the
corresponding receipt: for label $6$ the quadric recomputed from
coefficient orders $\le 22$ alone agrees with the full computation
to $3\times10^{-92}$ and the invariants to $7\times10^{-91}$,
consistent with (and now implied by) Lemma~\ref{lem:jets}.
\end{remark}

\begin{remark}[the field is $K$, not $\QQ$]
Even for the $\GQ$-fixed label $6$ the marked invariants need not be
rational: the marked triple $(X,\phi,P)$ is a $K$-object, since $P$
lies over $\sqm$. The computation below finds them in $K$ exactly.
Complex conjugation ($\sigma\notin\GK$) sends the $\sqm$-marked data
of label $i$ to the $-\sqm$-marked data of label $\kap(i)$; the
resulting compatibility between labels $\{1,2\}$ and $\{5,7\}$ is a
numerical gate in \S\ref{sec:partition}, not an assumption.
\end{remark}

We use the six weight-$0$ ratios
\[
  J_1=\frac{q_{13}}{q_{22}},\quad
  J_2=\frac{q_{11}q_{13}}{q_{12}^2},\quad
  J_3=\frac{q_{11}q_{33}}{q_{13}q_{22}},\quad
  J_4=\frac{q_{12}q_{23}}{q_{13}q_{22}},\quad
  J_5=\frac{q_{11}q_{23}}{q_{12}q_{13}},\quad
  J_6=\frac{q_{22}^2}{q_{11}q_{33}}.
\]
A numerical gate checks, for every label, that each denominator
factor appearing in $J_1,\dots,J_6$ is bounded away from zero
(receipts), so the six ratios are defined at all seven points.

\subsection{Ground-truth gate: label 6}
For label $6$, PSLQ at $100$ working digits recognizes all six
invariants in $K$ with small height:
\begin{center}
\LabelSixTable
\end{center}
Each recognition leaves $\sim 85$ digits of confirmation margin. As a
negative control, a ratio with mismatched weight (hence still
containing the transcendental $\lambda$) fails recognition, exactly
as Lemma~\ref{lem:weights} predicts.

% ============================================================ PARTITION
\section{The Galois orbit partition}\label{sec:partition}

For each invariant $J$ the multiset $\{J(X_1),\dots,J(X_7)\}$ (all
marked at the $\sqm$-side) is $\GK$-stable, so
\[
  \Phi_J(T) \;=\; \prod_{i=1}^{7}\bigl(T-J(X_i)\bigr)\;\in\;K[T],
\]
and the degrees of the irreducible factors of $\Phi_J$ over $K$
compute the $\GK$-orbit partition of the seven marked covers
(provided $J$ separates the seven, which is checked numerically).
The $\GQ$-partition then follows by gluing with
Theorem~\ref{thm:kappa}: $\GQ$-orbits are the $\kap$-stable unions of
$\GK$-orbits generated by the pairs $\{1,2\}$, $\{5,7\}$.

Gates: (i) agreement of the partition across the independent
invariants $J_1,\dots,J_6$; (ii) the $\kap$-consistency check that
the directly-computed invariants of labels $2$ and $7$ are the
appropriate Galois conjugates of those of $1$ and $5$; (iii) a
precision ladder (recognition at $60$ digits, confirmation at $100$;
escalation to $200$ digits if coefficient heights require it).

\subsection*{From partition to permutation group}
The factor degrees of $\Phi_J$ give only the orbit \emph{sizes}. From
the recognized (and consistency-checked) factors we extract more: the
Galois group of each irreducible factor over $K$, discriminants, and
Frobenius cycle types at a sample of good primes of $K$ determine the
image
\[
  A_K \;=\; \text{(the permutation group induced by }\GK\text{ on the
  seven)} \;\le\; S_7
\]
up to conjugacy, and adjoining the unconditionally known
complex-conjugation permutation gives
$A_\QQ=\langle A_K,\kap\rangle=\rho(G_\QQ)$ as an \emph{explicit
subgroup of $S_7$}. It is this group, not merely the partition
shape, that is to be placed alongside the ambient bound of
\S\ref{sec:gt}.

\subsection*{Epistemic status and the exactification path}
We state the outcome as a \emph{high-precision numerical
determination}: by Proposition~\ref{prop:equivariance} the quantities
$J(X_i)$ are honest algebraic numbers and the set-level Galois
structure is as described, but their \emph{identification} with
specific elements of $\overline{\QQ}$ rests on integer-relation
detection (PSLQ margin plus cross-implementation agreement). This is
compelling evidence, not yet proof. Two upgrades convert it to a
theorem. \emph{Cheap consistency}: the recognized values must satisfy
exact algebraic identities --- the recognized $\Phi_J$ must have
$K$-coefficients matching the recognized roots symmetrically, and
distinct invariants satisfy resultant compatibilities --- all
checkable in exact arithmetic; failures would falsify recognition.
\emph{Definitive}: reconstruct from the recognized invariants exact
models $(Q_i, C_i)$ over the recognized fields by lattice reduction,
and verify \emph{algebraically} that each defines a genus-$4$ curve
whose associated cover has monodromy $\Mtw$ of type
$(2A,23A,23B)$ --- at which point the partition statement inherits a
full proof. The first upgrade is carried out here: the recognized
sextic coefficients are the exact symmetric functions of the six
recognized roots by construction, and the two independently
recognized sextics satisfy an exact compatibility, verified purely
symbolically --- $f_2$ acquires a linear factor over
$\QQ[x]/(f_1)$, so the two independent descended coordinates
generate the \emph{same} sextic field $F$
(receipt \texttt{partition-p100-receipt.json}). The second upgrade
is future work in the spirit of \cite{HJLPPZ}'s own exactness step.

\subsection*{Descent to $\QQ$: trace and norm along the marking}
Let $A$ denote the degree-$7$ \'etale $\QQ$-algebra of the unmarked
Hurwitz fiber; the $\sqm$-marked fiber has coordinate algebra
$A_K=A\otimes_\QQ K$ with $K/\QQ$-involution $\tau$, and
Proposition~\ref{prop:equivariance} realizes each marked invariant as
an element $J\in A_K$. The elements
\[
  S=J+\tau J,\qquad N=J\cdot\tau J
\]
are the quadratic trace and norm and lie in $A$: they are
$\GQ$-equivariant functions of the \emph{unmarked} points. Pointwise,
by Theorem~\ref{thm:kappa}, $\tau$ acts by
$(\tau J)(i)=\overline{J(\kap(i))}$, so $S$ and $N$ are computable
from the $\sqm$-marked data alone. (The assembled polynomials
$\prod_i(T-S(i))$ have real coefficients \emph{by construction} ---
receipt \texttt{DESCENT\_ASSEMBLY\_REALITY}, an implementation gate,
not an equivariance test.)

\subsection*{The result}
Among small integer combinations of the twelve descended coordinates
$S_r$, $N_r$ ($r=1,\dots,6$), the coordinate $U_1=S_1+N_1$ (and,
independently, $U_2=S_2+N_2$) separates the seven points and
produces exact recognitions with $\sim70$ digits of margin at working
precision $100$:

\begin{numresult}[high-precision determination of the orbit partition]
\label{thm:partition}
Subject to the integer-relation recognitions below (margins
$10^{-88}$--$10^{-90}$, all gates passed), the $\GQ$-set of the seven
Hurwitz points is
\[
  \{6\}\;\sqcup\;\{1,2,3,4,5,7\},
\]
i.e.\ the partition is $1+6$: label $6$ is the unique rational point
and the remaining six form a single $\GQ$-orbit. Moreover the
associated permutation images are
\[
  \rho(\GQ)\;=\;\operatorname{Stab}_{S_7}(6)\;\cong\;S_6,
  \qquad
  \rho(\GK)\;=\;S_6 \ \text{as well.}
\]
Equivalently, $H\simeq\operatorname{Spec}\QQ\sqcup\operatorname{Spec}F$
with $F=\QQ[T]/(f)$ sextic, $\operatorname{Gal}(f^{\mathrm{spl}}/\QQ)
= S_6$, and signature $\operatorname{sig}(F)=(2,2)$ \textup{(}forced
by Theorem~\ref{thm:kappa}\textup{)}.
\end{numresult}

For the remainder of \S\S\ref{sec:partition}--\ref{sec:gt},
assertions identifying $F$ with the residue field of the six-point
Hurwitz orbit are conditional on the integer-relation recognitions
of Numerical Result~\ref{thm:partition}; statements about the
displayed polynomials $f$, $f_2$ themselves are exact.

\noindent Evidence and gates:
\begin{enumerate}
\item $U_1(6)=-\tfrac{5541}{10237}$ exactly (equal to
  $S_1(6)+N_1(6)$ from the label-$6$ ground truth of
  \S\ref{sec:invariants}), and the six values
  $U_1(i)$, $i\neq6$, are the roots of the exact monic sextic with
  common denominator $50019668646075143$:
  \begin{center}\small\SexticTable\end{center}
  Each of the six independently computed roots satisfies $f$ to
  $6\times10^{-88}$.
\item $f$ is irreducible over $\QQ$ with
  $\operatorname{Gal}=S_6$ (6T16), verified independently by
  \textsc{SymPy}, by a Frobenius inventory (below), and by a Magma
  referee run (receipt \path{magma-sextic-referee-ledger.txt}).
\item Replication and field identity: $U_2=S_2+N_2$ independently
  yields a second exact irreducible sextic $f_2$ with Galois closure
  of group $S_6$ and $U_2(6)=\tfrac{54447}{142129}$, consistent with
  the label-$6$ anchors. Moreover --- purely symbolically, with no
  numerical input --- $f_2$ acquires a linear factor over
  $\QQ[x]/(f_1)$: the two independently descended coordinates
  generate the \emph{same} sextic field $F$. This compatibility is
  an exact statement about the two displayed polynomials,
  independent of any numerics; the explicit isomorphism certificate
  is displayed in Appendix~\ref{app:beta}.
\item Frobenius inventory over $199$ good primes (the first $199$
  primes coprime to the sextic denominator and to the discriminant
  numerator, excluding $23$): nine distinct cycle types observed,
  with frequencies consistent with $S_6$ Chebotarev densities (the
  two absent classes are those of smallest density). Over the $94$
  \emph{split} primes $\bigl(\tfrac{-23}{p}\bigr)=+1$, whose
  Frobenii sample $\rho(\GK)$, the \emph{odd} classes $(1,2,3)$,
  $(6)$, $(1,1,4)$ occur $52$ times. This inventory serves as
  replication; the exact statement is Proposition~\ref{prop:gk}
  below.
\item The commutator-class experiment of
  Remark~\ref{rem:commutator} is settled: the six-point orbit crosses
  the partition $\{1,2,3,4\}\,|\,\{5,6,7\}$, so the arithmetic action
  does not stabilize it.
\end{enumerate}

\begin{proposition}[exact form of $\rho(\GK)=S_6$]\label{prop:gk}
Assume the recognitions of Numerical
Result~\ref{thm:partition}, and let $L$ be the splitting field of
$f$, so $\operatorname{Gal}(L/\QQ)=S_6$. Then $L\cap K=\QQ$ and
$\operatorname{Gal}(LK/K)\cong S_6$; consequently
$\rho(\GK)=S_6$.
\end{proposition}

\begin{proof}
An $S_6$-extension of $\QQ$ has a unique quadratic subfield, the
fixed field of $A_6$, namely
$\QQ(\sqrt{\operatorname{disc}f})$. Complex conjugation acts on the
six roots with cycle type $2^21^2$ (Theorem~\ref{thm:kappa}), an
\emph{even} permutation; it therefore lies in $A_6$ and fixes the
quadratic subfield pointwise, so that subfield is real
(equivalently $\operatorname{disc}f>0$, as the signature $(2,2)$
forces). Since $K$ is imaginary quadratic, $K\not\subset L$, hence
$L\cap K=\QQ$ and restriction is an isomorphism
$\operatorname{Gal}(LK/K)\xrightarrow{\;\sim\;}
\operatorname{Gal}(L/\QQ)=S_6$. The action of $\GK$ on the six
points is the action of $\operatorname{Gal}(LK/K)$ on the roots of
$f$, which is therefore all of $S_6$.
\end{proof}

\begin{remark}[the compositum explains the degree-$12$ wall]
$F$ has two real embeddings, hence contains no imaginary quadratic
field; so $F\cap K=\QQ$ and $[FK:\QQ]=12$, $[FK:K]=6$. The
$\sqm$-marked invariants of the six non-rational points naturally
generate subfields of $FK$: the degree-$12$ heights first encountered
by the marked computation were the compositum announcing itself.
\end{remark}

\subsection*{Local structure at $23$}
The rescaling $t\mapsto t/\sqm$, defined over $K$, moves the branch
locus to $(0,1,-1)$: the apparent collision of $0,\pm\sqm$ modulo
the prime above $23$ is a coordinate artifact, not an intrinsic
degeneration of the configuration. (No good-reduction claim at
other primes is intended: in the rescaled model $1$ and $-1$
collide modulo $2$, and $2$ divides the order of the $2A$
monodromy; no prime other than $23$ is analyzed here.) What singles
out $23$ is that it divides \emph{two} of the three local monodromy
orders --- the setting of stable reduction of three-point covers
with cyclic $p$-Sylow subgroup (Raynaud; Bouw--Wewers; Obus
\cite{ObusI,ObusII}). At $p=23$ the local decomposition of the
Hurwitz algebra is strikingly asymmetric.

\begin{proposition}[local factorization at $23$]\label{prop:local23}
$F\otimes\QQ_{23}\simeq F_2\times F_4$ with
$(e,f)=(2,1)$ and $(4,1)$ respectively; both extensions are totally
ramified and tame. Hence the seven-point Hurwitz algebra decomposes
locally as
\[
  \QQ_{23}\;\times\;F_2\;\times\;F_4 ,
\]
tame inertia at $23$ acts on the seven points with cycle type
$(1)(2)(4)$, and \emph{label $6$ is the unique inertia-fixed point of
the fiber at $23$}.
\end{proposition}

\begin{proof}[Proof (certificate)]
$f\bmod 23=(x+13)^2(x+9)^4$. The quartic cluster has a single Newton
segment of slope $-1/4$ (Eisenstein--Dumas: a totally ramified
quartic); the quadratic cluster has odd $23$-adic discriminant
valuation ($v_{23}(\operatorname{disc}f)=6=3+3$ with vanishing
cross-resultant valuation, the clusters lying at distinct residues),
hence is ramified. Both indices are prime to $23$ (tame). The
factorization, ramification indices, and residue degrees are
independently confirmed by a Magma computation over
$\QQ_{23}$ (\texttt{magma-local23-referee-ledger.txt}).
\end{proof}

\begin{remark}[towards a mechanism]
We emphasize the careful reading: Proposition~\ref{prop:local23}
exhibits label $6$ as the unique \emph{inertia}-fixed point at
$23$ --- a local distinction, weaker than global rationality, but
the first structural property found that singles out label $6$. Two remarks situate it in the stable-reduction
theory for cyclic $p$-Sylow (here $P=C_{23}$ of order exactly $23$,
$N_{\Mtw}(P)\cong 23{:}11$, so $m=|N/C|=11$). First, the proven
theorems of \cite{ObusI,ObusII} do not cover $\Mtw$: part~I assumes
$G$ $p$-solvable, part~II assumes $m=2$. However, for $|P|=p$ their
conclusion pattern --- vanishing of the first higher ramification
groups, i.e.\ \emph{tameness} of the field of moduli at $p$ --- is
exactly what Proposition~\ref{prop:local23} establishes for
$\Mtw$ with $m=11$: our data is a new data point consistent with
the expected phenomenology outside the proven range. Second, the
sharply posed open question: \emph{derive the inertia cycle type
$(1)(2)(4)$ --- and the identity of the fixed point --- from the
deformation data of the stable model at $23$, without using $f$}.
An affirmative answer would explain why this point is locally
distinguished; the remaining gap to global rationality would then be
the passage from the decomposition group at $23$ to all of $\GQ$.
\end{remark}

\PtwoHundredLine

% ================================================================== GT
\section{Finite Grothendieck--Teichm\"uller constraints}\label{sec:gt}

The computations above determine \emph{what} the Galois set of seven
points is; the question of \cite{HJLPPZ} --- \emph{why} is there a
fixed point --- asks for a structural cause. In this section we prove
a sharp negative result: the finite-level Grothendieck--Teichm\"uller
constraints attached to $\Mtw$ impose \emph{no} restriction at all on
permutations of the seven. Whatever forces the fixed point is
invisible at this level.

For a finite group $G$, Guillot \cite{Guillot} defines groups
$\mathcal{GT}(G)\supseteq\mathcal{GT}_1(G)$ acting on the set
$\mathscr{P}_c(G)$ of generating pairs of $G$ up to simultaneous
conjugacy, compatibly with the tower defining Drinfeld's
$\widehat{\mathcal{GT}}\supseteq G_\QQ$; the operators
$\theta\cdot[g,h]_c=[h,g]_c$ and
$\delta\cdot[g,h]_c=[h^{-1}g^{-1},h]_c$ generate an
$S_3$ of tautological symmetries. For $G$ simple and non-abelian he
proves the remarkable identification
\[
  \mathcal{GT}_1(G)\;\cong\;\mathscr{S}(G)
  \;:=\;C_{\operatorname{Sym}(\mathscr{P}_c)}\!\bigl(\operatorname{Out}(G)
  \times\langle\theta,\delta\rangle\bigr)\;\cap\;Y,
\]
where $Y$ is the subgroup preserving the partition of
$\mathscr{P}_c$ by ordered class-vectors, and moreover
\cite[Prop.~4.12]{Guillot} that $\mathscr{S}(G)$ is a product of
wreath products $E_k\wr S_{r_k}$, one factor per equivalence class of
$\langle\theta,\delta\rangle\times\operatorname{Out}(G)$-orbits, with
$E_k$ a subquotient of the acting group.

\begin{theorem}[no constraint at the $\Mtw$-level]\label{thm:gt}
The image of $\mathcal{GT}_1(\Mtw)$ in the symmetric group of the
set $\Ni(2A,23A,23B)^{\mathrm{in}}$ is the full $S_7$.
\end{theorem}

\begin{proof}
$\operatorname{Out}(\Mtw)=1$, so the acting group is
$H=\langle\theta,\delta\rangle\cong S_3$. The six ordered
class-vectors in the $H$-orbit of $(2A,23A)$ are pairwise distinct
(inversion exchanges $23A$ and $23B$ since $-1$ is a quadratic
non-residue mod $23$), so $H$ acts freely on the $42$ pair-classes
lying over them, with seven orbits --- one for each Nielsen class ---
and every orbit meets each of the six class-vector blocks in exactly
one point. These seven $H$-orbits are therefore equivalent in the
sense of \cite[Prop.~4.12]{Guillot}, and the corresponding wreath
factor of $\mathscr{S}(\Mtw)$ is $E\wr S_7$, where $E$ consists of
the $H$-equivariant self-maps of one free orbit preserving each block
singleton. An $H$-equivariant self-map of a free orbit is right
translation by some $s\in H$, which preserves the six singletons iff
$s=1$; hence $E=1$ and the factor is $S_7$, acting by arbitrary
permutation of the seven orbits. Since
$\mathcal{GT}_1(\Mtw)\cong\mathscr{S}(\Mtw)$, the restriction to the
seven Nielsen classes is all of $S_7$.
\end{proof}

\subsection*{Scope}
Theorem~\ref{thm:gt} concerns the \emph{single finite level}
$\mathcal{GT}_1(\Mtw)$. The comparison with the global
Grothendieck--Teichm\"uller action is inverse-limit in nature:
information at different finite levels must satisfy simultaneous
compatibility conditions, and cross-level compatibilities (for
instance those
exploited by Belyi-extending constructions, which compare the
$\Mtw$-projection with projections at other monodromy groups) are
\emph{not} eliminated by Theorem~\ref{thm:gt}. What the theorem does
show is that every invariant internal to the single level ---
every quantity attached functorially to
$\mathscr{P}_c(\Mtw)$ alone --- is constant on the seven. Our
finite-depth searches through cross-level constructions
(postcomposition by Belyi-extending maps of outer degree $\le 12$
with all recolorings; square-root and pullback probes) found no
constraint either, but that remains an empirical statement about
finite search depth, not a structural one.

\begin{remark}[cyclotomic remark, in its correct strength]
$K=\QQ(\sqm)$ is the quadratic subfield of $\QQ(\zeta_{23})$, and
$23A$, $23B$ are exchanged exactly by non-residue powering. For
$\sigma\in\GK$ the mod-$23$ cyclotomic character lies in
$(\mathbb{F}_{23}^\times)^2\cong C_{11}$, and since $11\nmid 7!$,
\emph{no nontrivial permutation of the seven can factor solely
through the quadratic-residue cyclotomic quotient}. This does not
decompose a Galois element into a cyclotomic part and an
$f_\sigma$-part --- the permutation induced by $(\lambda,f)$ depends
on both jointly --- and we do not claim the ``cyclotomic part'' of
the action is trivial in any stronger sense.
\end{remark}

\begin{remark}[a separating, non-GT invariant]\label{rem:commutator}
The class of the commutator $[x,y]$ \emph{does} separate the seven:
labels $1,2,3,4$ have $[x,y]$ of order $4$ (cycle type
$4^42^21^3$) and labels $5,6,7$ of order $6$ (type $6^23^22^21$),
uniformly over all six recolorings. By Theorem~\ref{thm:gt} this
partition \emph{cannot} be $\mathcal{GT}_1(\Mtw)$-invariant --- an
instructive example of how natural generating-pair invariants fail
to be Galois-functorial. The computation of
\S\ref{sec:partition} settles the question negatively: the
six-point $\GQ$-orbit crosses $\{1,2,3,4\}\,|\,\{5,6,7\}$, so the
arithmetic action does \emph{not} stabilize this partition (it is
merely $\kap$-stable).
\end{remark}

\subsection*{The sharpened question}
The comparison with the inverse limit requires care about which
Galois group enters: the tower of level-one groups
$\mathcal{GT}_1(G)$ receives not $G_\QQ$ itself but its derived
subgroup $G_\QQ'$ \cite{Guillot}. Since
$\rho(G_\QQ)\cong S_6$ surjectively onto the six-point orbit
(Numerical Result~\ref{thm:partition}), the image of the derived
subgroup is $\rho(G_\QQ')=[S_6,S_6]=A_6$, and the correct chain is
\[
  A_6\;=\;\rho(G_\QQ')
  \;\subseteq\;
  \operatorname{im}\bigl(\widehat{\mathcal{GT}}_1\to
  \mathcal{GT}_1(\Mtw)\bigr)\Big|_{H}
  \;\subseteq\;
  \mathcal{GT}_1(\Mtw)\Big|_{H}\;=\;S_7 ,
\]
in which Theorem~\ref{thm:gt} identifies the outermost group and
eliminates it as an explanation: \emph{the existence of the rational
point is not forced by the $\Mtw$-level Grothendieck--Teichm\"uller
constraints}. The single finite level places no restriction on the
seven points. The image of the global inverse-limit group is
potentially smaller --- any of $A_6$, $S_6$, $A_7$, $S_7$ is
compatible with the chain --- but determining it requires
cross-level compatibility information not computed here.

Numerical Result~\ref{thm:partition} determines the arithmetic image
itself, and the answer is as sharp as it could be:
\[
  A_{\mathrm{arith}}\;=\;\rho(G_\QQ)\;=\;
  \operatorname{Stab}_{S_7}(6)\;\cong\;S_6 ,
\]
\emph{the largest subgroup of $S_7$ compatible with fixing label
$6$} --- and already $\rho(\GK)=S_6$. The arithmetic action is
maximally mixing everywhere the fixed point permits: the action on
the six-point orbit is primitive, so it admits no nontrivial
invariant partition, and the quadratic field $K$ does not appear in
the Galois closure of the sextic
(Proposition~\ref{prop:gk}). The question of \cite{HJLPPZ} is
thereby concentrated into a single sharp statement: \emph{the entire
symmetry breaking is
\[
  S_7 \rightsquigarrow \operatorname{Stab}_{S_7}(6),
\]
one rational point, and otherwise maximal Galois mixing. Why does
arithmetic single out label $6$ --- and nothing else?}

% ============================================================ ARTIFACTS
\section{Verification and artifacts}\label{sec:artifacts}

All computations, receipts, and the Lean development are contained in
the artifact repository
\url{https://github.com/azide0x37/m23-seven-points}; releases of the
repository are archived with a DOI.
The receipt files are JSON with self-describing gates; each
computational claim above names its receipt. Independent
implementations: (a) a Python enumeration (bytes-encoded permutation
BFS), (b) a from-scratch re-implementation used by the Lean
certificate generator, (c) a Magma referee script (census,
generation, and $\kap$-certificates), run on the public Magma
calculator. An in-progress Lean~4 development encodes and checks
the small finite certificates (kernel-checked \texttt{decide}
proofs, with a single quarantined \texttt{native\_decide}
enumeration); the full census build and the glue lemmas connecting
those certificates to the paper's theorem statements are
incomplete, and no end-to-end formal verification is claimed here.

The analytic layer is validated as described in \S\ref{sec:powser};
we emphasize that the final arbiter for every recognized algebraic
number is integer-relation margin plus cross-implementation
agreement. \PtwoHundredArtifactLine

\subsection*{AI assistance and verification}
This project made extensive use of AI systems, principally Claude
(Anthropic) and ChatGPT (OpenAI), throughout the research process.
Their roles included exploratory mathematical discussion, proposing
and criticizing candidate approaches, literature triage, generation
and debugging of computational code, orchestration and analysis of
numerical experiments, and drafting and editing the manuscript under
the author's direction. The author selected the research questions,
determined which computational and mathematical claims to pursue,
executed or supervised the computations, evaluated the resulting
evidence, and takes responsibility for all claims in the paper.
Computational claims are accompanied by reproducible artifacts and
independent checks as described in this section; theoretical claims
are supported by the proofs and references given in the text.
AI-generated suggestions were not treated as mathematical evidence.

\subsection*{Acknowledgments}
The author thanks Bjorn Poonen for correspondence concerning
\cite{HJLPPZ}, including confirmation that the authors independently
computed the action of complex conjugation.
The author also thanks John A.~Janik for emphasizing, in an
October 5, 2026 discussion, the parity observation and the
distinction between forced real existence, the computed count,
and arithmetic rationality.

% ============================================================ APPENDIX
\appendix

\section{Representatives and certificates}\label{app:reps}

\begin{center}
\small
\GeneratorsTable
\end{center}

\medskip
Anchored representatives (labels of Convention~\ref{conv:group},
with label $6$ the class of the triple printed in
\cite{HJLPPZ}); every Nielsen
class is represented by the pair $(g_1,y_i)$:
\begin{center}
\small
\RepsTable
\end{center}

\medskip
\CGensLine.

\medskip
$\kap$-certificates ($\sigma=(1\,2)(5\,7)$):
\begin{center}
\small
\KappaCertsTable
\end{center}

\section{Computation parameters}\label{app:params}

\begin{center}
\ParamsTable
\end{center}

\section{The field-identity certificate}\label{app:beta}

Let $\alpha$ be a root of the sextic $f=f_1$ displayed in Numerical
Result~\ref{thm:partition}, and let $f_2$ be the second descended
sextic:
\begin{center}\small\SexticTwoTable\end{center}

\medskip
The element $\beta=\sum_{k=0}^{5} r_k\,\alpha^k\in\QQ[\alpha]/(f_1)$
with the coefficients below satisfies $f_2(\beta)=0$. The
verification is the exact polynomial-remainder identity
$f_2(g(x))\equiv 0 \pmod{f_1(x)}$ for $g(x)=\sum_k r_kx^k$ in
$\QQ[x]$ (receipt \texttt{iso-beta-receipt.json}); hence the two
independently recognized sextics define the same field $F$.
\begin{center}\small\IsoBetaTable\end{center}

\begin{thebibliography}{HJLPPZ}

\bibitem[HJLPPZ]{HJLPPZ}
X.~Huang, B.~Jackson, K.-H.~Lee, B.~Poonen, R.~Pries, S.~Zhang,
\emph{The Mathieu group $M_{23}$ is a Galois group over
$\mathbb{Q}$}, preprint (2026), arXiv:2608.08538.

\bibitem[Fri77]{Fried77}
M.~Fried, \emph{Fields of definition of function fields and Hurwitz
families --- groups as Galois groups}, Comm. Algebra \textbf{5}
(1977), 17--82.

\bibitem[DD97]{DebesDouai}
P.~D\`ebes, J.-C.~Douai, \emph{Algebraic covers: field of moduli
versus field of definition}, Ann. Sci. \'Ecole Norm. Sup. (4)
\textbf{30} (1997), 303--338.

\bibitem[KMSV14]{KMSV}
M.~Klug, M.~Musty, S.~Schiavone, J.~Voight, \emph{Numerical
calculation of three-point branched covers of the projective line},
LMS J. Comput. Math. \textbf{17} (2014), 379--430, arXiv:1311.2081.

\bibitem[VW14]{VoightWillis}
J.~Voight, J.~Willis, \emph{Computing power series expansions of
modular forms}, in \emph{Computations with modular forms}, Contrib.
Math. Comput. Sci. \textbf{6}, Springer (2014), 331--361,
arXiv:1205.0045.

\bibitem[Bel]{BelyiRepo}
M.~Musty et al., \emph{Belyi maps database and computation package},
\url{https://github.com/michaelmusty/Belyi} (MIT license); standalone
solver \texttt{Cext/powser\_arnoldi.c}.

\bibitem[Gui]{Guillot}
P.~Guillot, \emph{The Grothendieck--Teichm\"uller group of a finite
group and $G$-dessins d'enfants}, in \emph{Symmetries in graphs,
maps and polytopes}, Springer Proc.\ Math.\ Stat.\ \textbf{159}
(2016), 159--191, DOI 10.1007/978-3-319-30451-9\_8,
arXiv:1407.3112.

\bibitem[Ob11a]{ObusI}
A.~Obus, \emph{Fields of moduli of three-point $G$-covers with
cyclic $p$-Sylow, I}, Algebra Number Theory \textbf{6} (2012),
no.~5, 833--883, arXiv:0911.1103.

\bibitem[Ob11b]{ObusII}
A.~Obus, \emph{Fields of moduli of three-point $G$-covers with
cyclic $p$-Sylow, II}, J.\ Th\'eor.\ Nombres Bordeaux \textbf{25}
(2013), no.~3, 579--633, arXiv:1001.3723.

\bibitem[ATLAS]{ATLAS}
J.~H.~Conway, R.~T.~Curtis, S.~P.~Norton, R.~A.~Parker,
R.~A.~Wilson, \emph{Atlas of Finite Groups}, Oxford Univ.\ Press,
1985. (For $\operatorname{Out}(\Mtw)=1$ and the power map on the
classes $23A$, $23B$.)

\end{thebibliography}

\end{document}
